% CO1b (v1) -- Calculus of motion, PHYS 201, Fall 2024.  ANSWER KEY.
%
% Converted from "CO1b - Calculus Motion - v1 - Fall 2024 - Physics 201.docx",
% which is kept alongside this file as the source of record for the wording.
%
% This is the WORKED-SOLUTION version.  Its companion,
% CO1b-Fall2024-PHYS201-v1.tex, generates the same quiz blank for students.
% The question code below is identical to that file's; keep the two in step.
%
% Two separate records of the answers come out of this document:
%   * the worked steps printed into the PDF, for handing back or marking from;
%   * the randassign key, written to a separate file as a side effect of the
%     pythontex pass and never typeset here (see CLAUDE.md).
%
% The shared preamble lives in jjc-assessment.sty.
%
% To print a matching pair, set SEED to the same integer in both files (see the
% note in the pycode block below), then build each with the usual three passes:
%     pdflatex -interaction=nonstopmode <file>.tex
%     pythontex <file>.tex
%     pdflatex -interaction=nonstopmode <file>.tex

\documentclass[12pt]{exam}
\usepackage{jjc-assessment}

\renewcommand{\class}{Physics 201}
\renewcommand{\term}{Fall 2024}
\renewcommand{\examnum}{CO1b (v1)}

% jjc-assessment's running head rule sits close enough to the text block
% that a question starting at the top of a page touches it.  Set here rather
% than in the package, which other assessments share.
\setlength{\headsep}{0.3in}

% The Word original carries the work note as a running footer, which
% jjc-assessment does not define.  Only the footer slots are set up here: the
% \pagestyle switch has to come *after* \assessmentheader, because that macro
% selects \pagestyle{head} itself and would otherwise drop the footer again.
\firstpagefooter{}{\textbf{SHOW ALL WORK FOR HIGHEST RATING}}{}
\runningfooter{}{\textbf{SHOW ALL WORK FOR HIGHEST RATING}}{}

\begin{document}

\assessmentheader
\pagestyle{headandfoot}

% Marked on a line of its own rather than inside \examnum: the header
% tabular has fixed-width name and i.d. rules and overruns the text block
% if the exam number grows.
\begin{center}\textbf{\large ANSWER KEY}\end{center}

% The quiz body is generated from the single pycode block below, so every number
% on the page is computed rather than typed.  Question 1 is deliberately fixed;
% question 2 is drawn fresh for each paper.

\begin{pycode}
import random

import sympy as sp

from randassign import RandAssign
ra = RandAssign()

# ---------------------------------------------------------------------------
# Which paper this build produces
# ---------------------------------------------------------------------------
# Leave SEED as None for a fresh paper on every build.  Set it to an integer --
# the SAME integer here and in the student file -- to reproduce one exact paper,
# so the blank quiz and the worked key describe the same particle.
#
# Only the standard-library RNG is used, so seeding controls every draw in this
# document.  (`from pylab import *` would shadow `uniform` with the numpy
# generator, which random.seed() does not drive; see CLAUDE.md.)
SEED = None

paper_code = SEED if SEED is not None else random.randrange(1, 10**6)
random.seed(paper_code)

t = sp.Symbol('t', nonnegative=True)
tau = sp.Symbol('tau', nonnegative=True)


# ---------------------------------------------------------------------------
# Question 1 -- Artemis rocket.  Deliberately NOT randomised.
# ---------------------------------------------------------------------------
V_BURNOUT = 1800        # m/s reached at the end of the burn
T_BURN_MINUTES = 2      # "the first two minutes after liftoff"
MINUTE_WORDS = {1: 'one', 2: 'two', 3: 'three', 4: 'four', 5: 'five'}


def question_rocket():
    """Average acceleration over the burn, then the height reached.

    The rocket starts from rest and the acceleration is taken as constant, so
    v(t) = a t and the height is the integral of v across the burn.  Both are
    evaluated with sympy rather than by quoting a kinematic formula.
    """
    burn = T_BURN_MINUTES * 60                        # seconds
    a_avg = sp.Rational(V_BURNOUT, burn)              # (v - v0)/(t - t0), v0 = 0
    v_of_t = sp.integrate(a_avg, (tau, 0, t))         # = a t
    height = sp.integrate(v_of_t, (t, 0, burn))       # = a t^2 / 2

    word = MINUTE_WORDS[T_BURN_MINUTES]
    statement = (
        r'The Artemis rocket is estimated to be able to reach '
        r'$\SI{%.4g}{\meter\per\second}$ in the first %s minutes after liftoff.'
    ) % (V_BURNOUT, word)
    parts = [
        r'What is the average acceleration of the rocket during this time?',
        r'If air resistance were not a factor and if the rocket went straight '
        r'up, how high off the ground would it be at the end of the first %s '
        r'minutes of its flight?' % word,
    ]

    working = [
        (
            r'The rocket starts from rest, so the whole of '
            r'$\Delta v = \SI{%.4g}{\meter\per\second}$ is gained over '
            r'$\Delta t = \SI{%.4g}{\second}$:'
            r'\[ a_{\mathrm{avg}} = \frac{\Delta v}{\Delta t} '
            r'= \frac{\SI{%.4g}{\meter\per\second}}{\SI{%.4g}{\second}} '
            r'= \SI{%.3g}{\meter\per\second\squared} \]'
        ) % (V_BURNOUT, burn, V_BURNOUT, burn, float(a_avg)),
        (
            r'Holding that acceleration constant from rest gives '
            r'$v(t) = %s$, and the height is the area under it:'
            r'\[ h = \int_{0}^{\SI{%.4g}{\second}} v(t)\,\mathrm{d}t '
            r'= \tfrac{1}{2}\,a_{\mathrm{avg}}\,t^{2} '
            r'= \SI{%.6g}{\meter} = \SI{%.3g}{\kilo\meter} \]'
        ) % (sp.latex(v_of_t), burn, float(height), float(height) / 1000.0),
    ]

    key = [
        r'$a_{\mathrm{avg}} = \SI{%.3g}{\meter\per\second\squared}$' % float(a_avg),
        r'$h = \SI{%.6g}{\meter}$' % float(height),
    ]
    return statement, parts, working, key


# ---------------------------------------------------------------------------
# Question 2 -- particle on the x axis.  Randomised.
# ---------------------------------------------------------------------------
# x(t) = A t - B t^3, so v = A - 3B t^2 vanishes at t = sqrt(A/(3B)).  Drawing B
# and that instant, then setting A = 3 B t0^2, keeps the root a whole number of
# seconds the way the original paper's 24t - 2.0t^3 does.  The original draw
# (A = 24, B = 2, t0 = 2) is one of the possibilities.
def question_particle():
    b_halves = random.choice([2, 3, 4, 5, 6])         # B = 1.0 .. 3.0
    coeff_b = sp.Rational(b_halves, 2)
    t_zero = random.choice([2, 3])                    # instant when v = 0
    coeff_a = 3 * coeff_b * t_zero**2

    # Evaluate part (b) somewhere other than the instant found in part (a), so
    # the two parts cannot collapse onto the same answer.
    while True:
        t_eval = sp.Rational(random.choice([2, 3, 4, 5, 6]), 2)
        if t_eval != t_zero:
            break

    x_of_t = coeff_a * t - coeff_b * t**3
    v_of_t = sp.diff(x_of_t, t)
    a_of_t = sp.diff(v_of_t, t)

    roots = [r for r in sp.solve(sp.Eq(v_of_t, 0), t) if r.is_real and r >= 0]
    t_rest = max(roots)
    a_at_rest = a_of_t.subs(t, t_rest)
    a_at_eval = a_of_t.subs(t, t_eval)

    statement = (
        r'A particle moving along the $x$ axis has a position given by '
        r'$x = (%.4g\,t - %.1f\,t^{3})$ m, where $t$ is measured in s.'
    ) % (float(coeff_a), float(coeff_b))
    parts = [
        r'What is the magnitude of the acceleration of the particle at the '
        r'instant when its velocity is zero?',
        r'What is the acceleration of the particle at '
        r'$t = \SI{%.1f}{\second}$?' % float(t_eval),
    ]

    # Show the derivatives with decimal coefficients rather than the exact
    # rationals sympy carries internally, so the steps read the way the position
    # is printed on the paper.
    x_shown = sp.Float(float(coeff_a)) * t - sp.Float(float(coeff_b)) * t**3
    v_shown = sp.diff(x_shown, t)
    a_shown = sp.diff(v_shown, t)

    working = [
        (
            r'Differentiate the position twice:'
            r'\[ v(t) = \frac{\mathrm{d}x}{\mathrm{d}t} = %s, \qquad '
            r'a(t) = \frac{\mathrm{d}v}{\mathrm{d}t} = %s \]'
            r'$v = 0$ at the positive root $t = \SI{%.4g}{\second}$, so the '
            r'acceleration there is'
            r'\[ |a| = \SI{%.3g}{\meter\per\second\squared} \]'
        ) % (sp.latex(v_shown), sp.latex(a_shown), float(t_rest), float(abs(a_at_rest))),
        (
            r'The acceleration is linear in $t$, so evaluate $a(t) = %s$ '
            r'directly:'
            r'\[ a(\SI{%.1f}{\second}) = \SI{%.3g}{\meter\per\second\squared} \]'
            r'It is negative: the particle is slowing down while moving in the '
            r'$+x$ direction.'
        ) % (sp.latex(a_shown), float(t_eval), float(a_at_eval)),
    ]

    key = [
        r'$|a| = \SI{%.3g}{\meter\per\second\squared}$' % float(abs(a_at_rest)),
        r'$a = \SI{%.3g}{\meter\per\second\squared}$' % float(a_at_eval),
    ]
    return statement, parts, working, key


# ---------------------------------------------------------------------------
# Layout
# ---------------------------------------------------------------------------
SOLUTION_BOX = (
    r'\begin{mdframed}[backgroundcolor=black!5,linewidth=0.4pt,'
    r'innertopmargin=6pt,innerbottommargin=6pt]'
    r'\textbf{Solution.} %s'
    r'\end{mdframed}'
)


def emit(statement, parts, working):
    """One numbered question with its lettered parts, each worked through."""
    print(r'\question ' + statement)
    print(r'\begin{parts}')
    for part, steps in zip(parts, working):
        print(r'\part ' + part)
        print(SOLUTION_BOX % steps)
    print(r'\end{parts}')


rocket_statement, rocket_parts, rocket_working, rocket_key = question_rocket()
particle_statement, particle_parts, particle_working, particle_key = question_particle()

# The paper code identifies which draw this key belongs to; setting SEED to it
# in the student file reprints the matching blank quiz.
print(r'\hfill{\scriptsize Paper code: %d}\par' % paper_code)
print(r'\setcounter{question}{0}')
print(r'\begin{questions}')
emit(rocket_statement, rocket_parts, rocket_working)
print(r'\newpage')
emit(particle_statement, particle_parts, particle_working)
print(r'\end{questions}')

# One addsoln call per question, in the order the questions are printed, so the
# randassign key lines up question for question with the paper.
ra.addsoln(rocket_key)
ra.addsoln(particle_key)
\end{pycode}

\end{document}
